\documentclass{article}

\usepackage[margin=2cm]{geometry}
\usepackage{graphicx}
\usepackage{graphics}
\usepackage{hyperref}
\usepackage{psfrag}
\usepackage{ifthen}
\usepackage{array}
\usepackage{longtable}
\usepackage{color}
\usepackage{tikz}
\usepackage{amsmath, amsfonts, amsthm}
\usepackage[shortlabels]{enumitem}
\setcounter{secnumdepth}{0}
\newcommand{\N}{\mathbb{N}}
\newcommand{\R}{\mathbb{R}}

\usetikzlibrary{svg.path}


\begin{document}

\title{Homework 5}
\author{Hayden Pott}
\date{22 April 2024}

\maketitle

\section{2-20}

Prove part 2 of Theorem 2.36.

\begin{proof}
    \begin{align*}
        A \setminus (B \cap C) &= \{x : x \in A \land x \not\in (B \cap C) \} \\
        &= \{x : x \in A \land \neg(x \in (B \cap C)) \} \\
        &= \{x : x \in A \land \neg(x \in B \land x \in C) \} \\
        &= \{x : x \in A \land (x \not\in B \lor x \not\in C) \} \\
        &= \{x : (x \in A \land x \not\in B) \lor (x \in A \land x \not\in C) \} \\
        &= \{x : x \in (A \setminus B) \lor x \in (A \setminus C) \} \\
        &= \{x : x \in ((A \setminus B) \cup (A \setminus C)) \} \\
        &= (A \setminus B) \cup (A \setminus C)
    \end{align*}

\end{proof}

\section{2-24}

Alternate proofs for Theorems 2.35 and Theorem 2.36.  For each of the following
proofs, use Standard Proof Technique 2.26 (Proof by mutual containment)

\begin{enumerate}
    \item[b)] Prove part 2 of Theorem 2.35
        \begin{proof}
            Let $x \in A \cup (B \cap C)$ be ABF. \\
            $\implies$ $x \in A \lor x \in B \cap C$ \\
            $\implies$ $x \in A \lor (x \in B \land x \in C)$ \\
            $\implies$ $(x \in A \lor x \in B) \land (x \in A \lor x \in C)$ \\
            $\implies$ $x \in (A \cup B) \land x \in (A \cup C)$ \\
            $\implies$ $x \in ((A \cup B) \cap (A \cup C))$ \\
            $\implies$ $A \cup (B \cap C) \subseteq (A \cup B) \cap (A \cup C)$ \\

            Let $y \in (A \cup B) \cap (A \cup C)$ be ABF. \\
            $\implies$ $y \in (A \cup B) \land y \in (A \cup C)$ \\
            $\implies$ $(y \in A \lor y \in B) \land (y \in A \lor y \in C)$ \\
            $\implies$ $y \in A \lor (y \in B \land y \in C)$ \\
            $\implies$ $y \in A \lor y \in (B \cap C)$ \\
            $\implies$ $y \in A \cup (B \cap C)$ \\
            $\implies$ $(A \cup B) \cap (A \cup C) \subseteq A \cup (B \cap C)$ \\

            Therefore, by mutual containment, we have that $A \cup (B \cap C) =
            (A \cup B) \cap (A \cup C)$
        \end{proof}
\end{enumerate}

\section{2-32}

Let $A$, $B$, $C$ be sets.  Prove each of the following equalities using the
Axiom of Extension.  That is, prove each of the following equalities by proving
that each element of the set on the left side is an element of the set on the
right side and vice versa.

\begin{enumerate}
    \item[d)] $A \Delta (B \Delta C) = (A \Delta B) \Delta C$
        \begin{proof}
            \begin{align*}
                A \Delta (B \Delta C) &= (A \setminus (B \Delta C)) \cup ((B \Delta C) \setminus A) \\
                &= \Big(A \setminus ((B \setminus C) \cup (C \setminus B))\Big) \cup \Big(((B \setminus C) \cup (C \setminus B)) \setminus A\Big) \\
                &= \Big((A \setminus (B \setminus C) \cap A \setminus (C \setminus B))\Big) \cup \Big(((B \setminus C) \cup (C \setminus B)) \setminus A\Big) \\
                &= ((A \setminus (B \setminus C)) \cap (A \setminus (C \setminus B))) \cup ((B \setminus C)\setminus A) \cup ((C \setminus B)\setminus A) \\
                &= (((A \setminus B \setminus C) \cup (B \setminus A \setminus C))) \cup (((C \setminus (A \setminus B)) \cap (C \setminus (B \setminus A))) \\
                &= (((A \setminus B) \setminus C \cup (B \setminus A) \setminus C)) \cup (((C \setminus (A \setminus B)) \cap (C \setminus (B \setminus A))) \\
                &= (((A \setminus B) \cup (B \setminus A)) \setminus C) \cup (C \setminus ((A \setminus B) \cup (B \setminus A)) \\
                &= ((A \Delta B) \setminus C) \cup (C \setminus (A \Delta B)) \\
                &= (A \Delta B) \Delta C 
            \end{align*}
        \end{proof}
\end{enumerate}

\section{2-36}

Power sets.

\begin{enumerate}[a)]
    \item Compute the power set of the set $\{a, b, c, d\}$ \\
        \textbf{ANSWER:} \\
        $P(S) = \{
        \emptyset,
        \{a\},
        \{b\},
        \{c\},
        \{d\},
        \{a,b\},
        \{a,c\},
        \{a,d\},
        \{b,c\},
        \{b,d\},
        \{c,d\},
        \{a,b,c\},
        \{a,b,d\},
        \{a,c,d\},
        \{b,c,d\},
        \{a,b,c,d\}
    \}$
    \item Compute the power set of the set $\{1, 2, 3, 4\}$ \\
        \textbf{ANSWER:} \\
        $P(S) = \{
        \emptyset,
        \{1\},
        \{2\},
        \{3\},
        \{4\},
        \{1,2\},
        \{1,3\},
        \{1,4\},
        \{2,3\},
        \{2,4\},
        \{3,4\},
        \{1,2,3\},
        \{1,2,4\},
        \{1,3,4\},
        \{2,3,4\},
        \{1,2,3,4\}
    \}$
    \item Is there any significant difference between Exercises 2-36a and 2-36b? \\
        \textbf{ANSWER}: No, the only difference is what we call the objects in
        the set
\end{enumerate}

\section{2-39}

Let $A = \{1, 2, 3, 4\}$ and let $B = \{a, b, c, d, e\}$.  Determine which of
the following relations is a function.  If it is a function, determine if it is
injective, surjective, bijective, or none of the above.  Justify your answers
using the definitions.

\begin{enumerate}[a)]
    \item $\{(1, a), (2, b), (3, c), (4, d)\}$ from $A$ to $B$ \\
        \textbf{ANSWER:} \\
        Function -- Well defined and totally defined \\
        Injective \checkmark -- No two values from $A$ map to the same value in
        $B$. \\
        Surjective $\times$ -- The value $e$ in $B$ can not be reached in this
        function.
    \item $\{(1, a), (2, b), (3, c), (4, d)\}$ from $A \setminus \{4\}$ to $B$ \\
        \textbf{ANSWER:} \\
        Function -- Well defined and totally defined \\
        Injective \checkmark -- No two values from $A \setminus \{4\}$ map to the same value in
        $B$. \\
        Surjective $\times$ -- The values $e$ and $d$ in $B$ can not be reached in this
        function.
    \item $\{(1, a), (2, b), (3, c), (4, c)\}$ from $A$ to $B$ \\
        \textbf{ANSWER:} \\
        Function -- Well defined and totally defined \\
        Injective $\times$ -- The value $c$ in $B$ can be reached by both $3$
        and $4$ in $A$. \\
        Surjective $\times$ -- The values $e$ and $d$ in $B$ can not be reached in this
        function.
    \item $\{(1, a), (2, b), (3, c), (4, a)\}$ from $A$ to $B \setminus \{d, e\}$ \\
        \textbf{ANSWER:} \\
        Function -- Well defined and totally defined \\
        Injective $\times$ -- The value $1$ in $B$ can be reached by both $1$
        and $4$ in $A$. \\
        Surjective \checkmark -- All values in $B$ can be reached from a value
        in $A$.
\end{enumerate}

\section{2-40}

For each pair of functions, check if they are inverses of each other.  For this
exercise, you may use what you know about real numbers and functions of real
numbers.

\begin{enumerate}[a)]
    \item $f, g : \R \to \R, f(x) = 3x + 6, g(x) = \frac{x}{3} - 2$
        \begin{proof}
            \begin{align*}
                f(g(x)) &= 3(\frac{x}{3} - 2) + 6 \\
                        &= x - 6 + 6 \\
                        &= x \\
                        &= id_\R(x) \\
                \\
                g(f(x)) &= \frac{3x + 6}{3} - 2 \\
                        &= x + 2 - 2 \\
                        &= x \\
                        &= id_\R(x) \\
            \end{align*}
        \end{proof}
        Thus, because $f \circ g$ and $g \circ f$ are both the identity
        function, $id_\R$, we can conclude that $f$ and $g$ are inverses of
        each other.
\end{enumerate}

\section{2-42}

Determine which of the following functions are one-to-one.  For those that are
one-to-one, compute the inverse from the range to the domain.  For this
exercise, you may use what you know about real numbers and functions of real
numbers.

\begin{enumerate}
    \item[c)] $f: \R \setminus \left\{\frac{5}{2}\right\} \to \R, f(x) =
        \frac{5x+4}{2x-5}$ \\
        \textbf{ANSWER:} \\
        One-to-one function \checkmark \\
        \begin{align*}
            x &= \frac{5g(x)+4}{2g(x)-5} \\
            x(2g(x)-5) &= 5g(x)+4 \\
            2xg(x) - 5x &= 5g(x) + 4 \\
            2xg(x) - 5g(x) &= 5x + 4 \\
            (2x - 5)g(x) &= 5x + 4 \\
            g(x) &= \frac{5x + 4}{2x - 5}\\
        \end{align*}
        Inverse: $g(x) = \frac{5x + 4}{2x - 5}$
\end{enumerate}

\section{2-43}

Each of the graphs in Figure 2.13 depicts a relation from $\R$ to $\R$.  For
each graph, decide if the relation is an injective, surjective, or bijective
function.

\begin{enumerate}[1)]
    \item Injective $\times$ \\
          Surjective $\times$ \\
          Bijective $\times$ \\

    \item Injective \checkmark \\
          Surjective \checkmark \\
          Bijective \checkmark \\

    \item Injective $\times$ \\
          Surjective \checkmark \\
          Bijective $\times$ \\

    \item Not a function \\
          Injective $\times$ \\
          Surjective $\times$ \\
          Bijective $\times$ \\
\end{enumerate}

\end{document}
